\documentclass[12pt,a4paper]{report}
\usepackage[utf8]{inputenc}
\usepackage[T1]{fontenc}
\usepackage[italian]{babel}
\usepackage{geometry}
\geometry{margin=2.5cm}
\usepackage{setspace}
\onehalfspacing
\usepackage{hyperref}
\usepackage{fancyhdr}
\usepackage{csquotes}

\pagestyle{fancy}
\fancyhf{}
\fancyhead[L]{Titolo della Tesina}
\fancyhead[C]{Nome Cognome}
\fancyhead[R]{\thepage}

\title{TITOLO DELLA TESINA}
\author{Nome Cognome}
\date{Anno Accademico 2024–2025}

\begin{document}

\begin{titlepage}
\centering
{\huge \textbf{TITOLO DELLA TESINA} \par}
\vspace{0.5cm}
{\large Sottotitolo o descrizione breve dell'argomento trattato \par}
\vspace{1cm}
\textbf{Disciplina:} Nome della materia \par
\vspace{0.5cm}
\textbf{Candidato/a:} Nome Cognome \par
\textbf{Matricola:} XXXXXXXX \par
\vspace{0.5cm}
\textbf{Relatore/Docente:} Prof. Nome Cognome \par
\vspace{0.5cm}
\textbf{Anno Accademico / Scolastico} 2024–2025 \par
\end{titlepage}

\tableofcontents
\newpage

\section{Introduzione}
Scrivi qui il testo introduttivo. L'introduzione dovrebbe presentare l'argomento trattato, le motivazioni della scelta, gli obiettivi della tesina e una breve descrizione della struttura del lavoro.

\section{Primo Capitolo}

\subsection{Primo sottoparagrafo}
Inserisci qui il contenuto del primo sottoparagrafo. Il testo è giustificato e utilizza l'interlinea 1,5 per una migliore leggibilità.

\subsection{Secondo sottoparagrafo}
Inserisci qui il testo del secondo sottoparagrafo.

\subsubsection{Esempio di citazione}
\begin{quote}
``Le citazioni dirette di testo vanno formattate così, con il rientro laterale e la barra colorata. Indicare sempre la fonte.'' (Autore, Anno, p. XX)
\end{quote}

\subsection{Elenchi}
Esempio di elenco puntato:
\begin{itemize}
\item Primo elemento dell'elenco
\item Secondo elemento dell'elenco
\item Terzo elemento dell'elenco
\end{itemize}

Esempio di elenco numerato:
\begin{enumerate}
\item Prima voce numerata
\item Seconda voce numerata
\item Terza voce numerata
\end{enumerate}

\section{Secondo Capitolo}
Continua qui con il secondo capitolo della tua tesina. Puoi aggiungere tutti i capitoli e sottoparagrafi necessari seguendo la struttura proposta.

\section{Conclusioni}
Scrivi qui le conclusioni del tuo lavoro. Riassumi i punti principali trattati, le riflessioni personali e gli eventuali sviluppi futuri dell'argomento.

\section{Bibliografia e Sitografia}

\subsection{Libri e articoli}
\begin{enumerate}
\setcounter{enumi}{3}
\item Cognome, N. (Anno). \textit{Titolo del libro}. Editore.
\item Cognome, N. (Anno). ``Titolo articolo''. \textit{Nome Rivista}, vol. X, pp. 00–00.
\end{enumerate}

\subsection{Risorse web}
\begin{enumerate}
\setcounter{enumi}{5}
\item Autore (Anno). \textit{Titolo pagina}. URL: \url{https://www.esempio.it}
\end{enumerate}

\end{document}